\section*{Operational category definitions}
These definitions complement the taxonomy panels. Evidence provenance refers to individual extracted claims and may overlap; observation and repository origin are not mutually exclusive. Analysis mode concerns the dominant mechanism, not an entire publication. Separate static and dynamic mechanisms do not alone establish a hybrid mechanism. Unreported information, uncertainty, inapplicability, and a negative assessment remain distinct. Category descriptions do not establish effectiveness or general correctness.
\begingroup
\setlength{\tabcolsep}{4pt}
\begin{longtable}{@{}p{.15\linewidth}p{.27\linewidth}p{.25\linewidth}p{.27\linewidth}@{}}
\caption{Operational definitions, inclusion criteria, and interpretation limits.}\label{tab:operational-definitions}\\
\toprule Category & Definition & Inclusion & Exclusion and limits\\\midrule
\endfirsthead
\toprule Category & Definition & Inclusion & Exclusion and limits\\\midrule
\endhead
\bottomrule\endfoot
Hypothetical example & A constructed example used to explain a possible quality problem, without an observation claim for that example. & An explanatory quality claim presented as hypothetical or illustrative. & Do not treat an illustration as an observed defect or infer incidence. Execution reported elsewhere does not determine this claim's provenance.\\
Mutation operator & A transformation intended to produce a faulty variant for assessment. & A source-described rule for modifying an artifact to introduce a fault. & An operator definition alone establishes neither an instantiated fault nor mutation-testing effectiveness.\\
Real observed & A fault or quality problem reported as actually observed in the investigated software context. & A source-located claim of observation, retaining its reported context and confirmation limits. & Constructed examples and injections are not naturally observed defects. Observation and repository provenance may coexist.\\
Repository reported & A quality claim whose provenance includes a software repository or issue report. & A source-located claim derived from repository or issue-report evidence, whether or not subsequently reproduced. & Repository origin alone proves neither independent confirmation nor field prevalence. Subsequent observation does not erase that origin.\\
Synthetic injected & A fault deliberately inserted into an artifact for an investigation. & A source-described instantiated injection, distinguished from the rule that generates it. & A proposed injection recipe alone is an operator, not evidence of an injected instance or a naturally occurring defect.\\
Taxonomy only & A quality class described taxonomically without a separately supported observed instance for the extracted claim. & A taxonomic description whose evidential role is conceptual classification. & Do not characterize every claim in a taxonomy paper this way; separately evidenced examples retain their own provenance.\\
Unclear evidence provenance & Relevant evidence is present, but its provenance cannot be reliably determined. & An explicit uncertainty about the provenance of the extracted quality claim. & Do not interpret uncertainty as a negative result, an observation, missing reporting, or inapplicability.\\
Relational and property-based testing & Testing based on properties of executions or relations among inputs, transformations, outputs, or implementations. & Metamorphic, differential, or property-based testing, preserving their distinct comparison rules. & A relational description alone does not establish a statistical test or effectiveness; other mechanisms need their own evidence.\\
Mutation-based assessment & Assessment of tests or assessment mechanisms using deliberately modified faulty variants. & A mechanism reported as mutation testing. & An isolated mutation-operator description or injected fault does not alone establish a mutation-testing assessment.\\
Debugging and fault localization & Mechanisms that diagnose or locate faults rather than merely report an unexpected outcome. & Source-described debugging or fault-localization mechanisms. & Failure detection alone and generic repair do not establish a diagnostic or localization mechanism.\\
Model-based testing & Testing in which a model has an explicit role in specifying, deriving, or organizing tests. & A mechanism reported as model-based testing with an identified testing role for the model. & Calling a circuit or artifact a model is insufficient. Frequency of reported support is not a definitional requirement.\\
Dynamic analysis and monitoring & Inspection or monitoring of execution behavior for an identified analytical purpose. & Source-described dynamic-analysis or runtime-monitoring mechanisms. & Execution alone does not establish this family, a test level, or effectiveness; other testing mechanisms remain distinct.\\
Dynamic analysis mode & The dominant assessment mechanism depends on execution of the target program or circuit. & Assessment requiring target execution and its resulting behavior or observations. & Running a checker is not necessarily target execution. Do not infer effectiveness, deployment, or test level.\\
Formal analysis mode & The dominant mechanism uses proof or logical reasoning about a property under stated assumptions. & A source-supported proof or logical assessment role. & Formal notation alone is insufficient. No universal precedence over static or symbolic mechanisms or general correctness guarantee is implied.\\
Hybrid analysis mode & A single assessment mechanism explicitly combines static and dynamic analysis. & A source-located mechanism integrating both forms of analysis. & Separate static and dynamic mechanisms in one publication, or quantum--classical computation alone, do not establish this mode.\\
Static analysis mode & The dominant mechanism inspects an artifact without requiring execution of the target program or circuit. & An artifact-analysis mechanism supported by the source's static assessment description. & Absence of target execution does not impose precedence over a formal proof role; terminology alone does not resolve ambiguous cases.\\
Assertion & A source-explicit predicate used to assert a property of program or model state or behavior. & A condition identified as an assertion in its stated context. & Do not infer an executable checker from a formula or add an invariant classification solely because the predicate may persist.\\
Invariant & A property intended to hold over the states, steps, or interval specified by the source. & A source-explicit invariant with its stated scope and assumptions. & Do not extend a scoped invariant to all executions, require stronger persistence, or infer implemented enforcement.\\
Checkability: No & A reported negative judgment about executability or checkability within the assessed scope. & The source-scoped checkability judgment is negative. & Missing reporting and ambiguous mechanisms are distinct from a negative judgment. This is not an effectiveness assessment.\\
Checkability: Not reported (NR) & Executability or checkability is not reported sufficiently to characterize it. & The relevant checking information is absent or insufficiently reported. & Do not infer impossibility, absence of any implementation, or a negative checkability judgment.\\
Checkability: Unclear & Relevant checking evidence is present but does not support a reliable checkability judgment. & The relationship between a condition and its checking mechanism remains ambiguous. & Do not resolve the ambiguity from notation alone or treat it as negative or unreported evidence.\\
Checkability: Yes & The condition is reported as executable or checkable within the stated scope. & A positive source-scoped checkability judgment under its stated assumptions. & Checkability does not establish empirical evaluation, effectiveness, scalability, deployment, or universal validity.\\
\end{longtable}
\endgroup
